\subsection*{13.5 Conditional Expectation Given a Continuous Random Variable}

Variable

While the process of deriving the formula for computing conditional expectation in

the continuous case is similar to that in the discrete case, it is more complex due to

the presence of events with probability zero.

\begin{thmbox}{13.14}
\end{thmbox}

Suppose .(R2,\mathcal{B}(R2), P) is a probability space on R2, and X and Y denote

the x - and y -coordinates, respectively. Let g(x) denote a Borel function in

L1(P) .

Suppose the joint probability density function of X and Y is denoted by

fXY (x, y), ie., for any Borel set B in\mathcal{B}(R2), we have

P(B) =

\int

B

fXY (x,y)d\lambda(x,y),

where. \lambda is the Lebesgue measure on R2. Then, we have

E[g(X)\vertY]=

\int

R

g(x)fX\vertY (x\verty) d\lambda(x), (13.13)

where

fX\vertY (x\verty)\coloneqq fXY (x, y)

fY (y)

and

fY (y)\coloneqq

\int

R

fXY (x, y) d\lambda(x). (13.14)

In particular , if X is integrable, we can compute the conditional expecta-

tionE[X\vertY] by

E[X\vertY]=

\int

R

xfX\vertY (x, y) d\lambda(y).

We may assumefY (y)/=0 for simplicity so that we do not need to worry about

division by zero. We note that the conditional expectation E[g(X)\vertY] in (13.13) is

a\sigma(Y) -measurable function and hence is a function of y .

\textit{Proof} Let h(y)=

\int

R g(x)fX\vertY (x\verty)d\lambda(x) be the candidate function for the

conditional expectation. Our goal is to prove that, for any \sigma(Y) -measurable set B ,

the following equation holds:

\int

B

g(x)dP(x,y) =

\int

B

h(y) dP(x, y). (13.15)

Consider a special type of \sigma(Y) -measurable set of the form B= R \times[ a,b] for

some a<b Assuming we can apply Fubini's theorem to evaluate the integrals by

iterated integration, we will show that Eq. (13.15) holds for this type of set. When

B is of the form R\times[ a,b], we can express the right-hand side of (13.15) as

\int

[a,b]

\int

R

h(y)fXY (x, y) d\lambda(x)d\lambda(y)=

\int

[a,b]

h(y)

\int

R

fXY (x, y) d\lambda(x)d\lambda(y)

=

\int

[a,b]

h(y)fY (y) d\lambda(y)

=

\int

[a,b]

\int

R

g(x)fXY (x, y) d\lambda(x)d\lambda(y).

The last line is the same as the left-hand side of (13.15).

We can justify the application of Fubini's theorem by assuming that \vertg(x)\vert is

integrable, as

\int

R\times[a,b]

\vertg(x)\vertfXY (x,y)d\lambda(x,y) \leq

\int

R\timesR

\vertg(x)\vertfXY (x,y)d\lambda(x,y) < \infty .

We obtain the following equation for all closed intervals I in R:

E[g(X) 1R\timesI ]= E[h(Y) 1R\timesI ]. (13.16)

We can extend (13.16) to all Borel sets B\in \mathcal{B} (R) by appealing to the \pi-\lambda

theorem. First, we observe that the collection of all closed intervals in R form a

\pi-system P. Next, we consider the collection L consisting of subsets S of \Omega that

satisfies

E[g(X) 1R\timesS]= E[h(Y) 1R\timesS].

By what we have proved above, we know that P\subset L We can prove that L is

indeed a \lambda-system, which implies that L contains all Borel sets in R. Therefore,

(13.16) holds for all Borel sets B\in \mathcal{B} (R).

Consequently, we have E[g(X) 1C]= E[h(Y) 1C] for all \sigma(Y) -measurable

sets C This implies that h(Y) is indeed a version of the conditional expectation

E[g(X)\vertY] , as it satisfies the defining properties of the conditional expectation. \hfill $\square$

The second part of the theorem establishes the equivalence between the measure-

theoretic definition of conditional expectation and the definition typically presented

in a first course of probability for continuous random variables. By establishing this

equivalence, we demonstrate that the abstract definition of conditional expectation

generalizes the familiar notion from introductory probability theory.

While the classical approach to conditional expectation is well-suited for dealing

with continuous or discrete random variables, it is not capable of handling the case

where we wish to compute the conditional expectation of a discrete random variable

given a continuous random variable, or vice versa. In such cases, the measure-

theoretic definition provides a more general and rigorous framework for defining

conditional expectation.

\textbf{Example 13.5.1 (Conditional Probability Given a Continuous Random Variable)} \\

Consider a point .(X, Y) chosen uniformly at random in the interior of a unit square with vertices

(0,0),.(0,1),.(1,1),a n d.(1,0)Let Z be a binary random variable defined by

Z\coloneqq

{

1i f Y<X

0i f Y\geq X.

We want to find E[Z\vertX]Note that X is uniformly distributed between 0 and 1, while Z is a

Bernoulli random variable.

Let A be the event .{Y<X }The random variable Z is the same as the indicator function . 1A,

and E[Z\vertX] is the conditional probability P(A \vertX)From the geometry of the problem, we can

make a guess that P(A \vertX)= X Indeed, we will show that

\int

B

ZdP =

\int

B

XdP (13.17)

holds for all \sigma(X)-measurable sets B .

ConsiderB=[ a,b]\times[ 0,1],w h i c h i s\sigma(X)-measurable. We have

\int

[a,b]\times[0,1]

ZdP =

\int

[a,b]\times[0,1]

1A dP= b2 -a 2

2 .

This is the area of a trapezoid with vertices .(a,0),.(b,0),.(b, b),a n d.(a, a).

On the other hand, we have

\int

B

XdP =

\int b

a

\int 1

xd y dx=

\int b

a

xdx = b2 -a 2

2 .

Hence, (13.17) holds for all set B that can be written as a rectangle .[a,b ]\times[ 0,1] for somea<b .

Using the\pi-\lambda theorem, we can extend this to all sets B that can be written as a Cartesian product

B= C \times[ 0,1],w h e r e C is Borel set in R. Finally, we note that X is\sigma(X)-measurable. Therefore,

X is a version of the conditional expectation E[1A\vertX], which is equal to P(A \vertX) by the definition

of conditional expectation. This proves that E[Z\vertX]= P(A \vertX)= X .
